% !TeX root = ../TCSVT_template/manuscript.tex
% Replacement for subsection C: two displayed equations only.
% Requires amsmath and the bibliography keys from
% mmdit_backbone_concise_references.bib.
% Retains eq:mmdit_joint_attention for subsection D.
% Removed formula labels: eq:mmdit_input_streams, eq:mmdit_stream_update,
% eq:mmdit_speaker_condition, eq:mmdit_velocity_output.
% The architecture and conditioning pathways are unchanged.

\subsection{Multimodal Diffusion Transformer for Video Dubbing}
\label{subsec:mmdit_backbone}
\label{sec:mmdit_dubbing}

Video dubbing requires coordinated modeling of linguistic content, visual
articulation, and speaker characteristics. To accommodate their different
representations and temporal structures, we adopt a multimodal diffusion
transformer (MM-DiT)~\cite{esser2024scaling} with separate audio, visual,
and textual streams. The backbone combines multimodal joint attention with
subsequent audio-only refinement~\cite{cheng2025mmaudio,yun2026immersivetts},
allowing cross-modal interaction while retaining modality-specific processing.

\subsubsection{Multimodal Representations}
\label{subsubsec:mmdit_inputs}
The audio stream receives the channel-wise concatenation of the noisy latent
$\mathbf{z}_t$ and the reference-latent context. During training,
complementary temporal masking~\cite{choi2025aligndit} retains acoustic
context outside the generation region and visual features within it; visual
positions corresponding to the acoustic context are replaced with a null
embedding. At inference, the reference utterance supplies the acoustic
context without requiring its corresponding video.

Lip-motion features from a pretrained audiovisual speech
encoder~\cite{shi2022avhubert} are linearly resampled from 25~Hz to the
40-Hz audio-latent grid. The transcript is encoded at the character level
using four ConvNeXt V2 blocks~\cite{woo2023convnext,chen2025f5tts},
without duration-based expansion. Learned input mappings project all three
streams to a common hidden width, while the text retains its own sequence
length $L_y$. The audio and visual pathways additionally incorporate temporal
contextualization and convolutional positional embeddings.

\subsubsection{Multimodal Joint Attention}
\label{subsubsec:mmdit_joint}
Each joint block uses modality-specific normalization, attention projections,
and feed-forward networks. Flow-time-conditioned adaptive layer
normalization~\cite{peebles2023dit} modulates the streams before attention.
Their queries and keys are RMS-normalized, with temporal rotary
embeddings~\cite{su2024roformer} applied to the audio and visual streams
on the shared time grid. The queries, keys, and values from the three
modalities are concatenated along the sequence dimension to form
$\mathbf{Q}$, $\mathbf{K}$, and $\mathbf{V}$. Joint attention for one head
is computed as
\begin{equation}
\mathbf{O}=\operatorname{softmax}\!\left(
\frac{\mathbf{Q}\mathbf{K}^{\top}}{\sqrt{d_h}}+\mathbf{B}
\right)\mathbf{V},
\label{eq:mmdit_joint_attention}
\end{equation}
where $d_h$ is the head dimension and $\mathbf{B}$ is the audiovisual
temporal bias defined in the next subsection; padding masks are omitted
for clarity. A single softmax normalizes over keys from all modalities.
The output is split back into the three streams, each of which is updated
through its own output projection, gated residual connections, and
feed-forward network. Thus, audio, visual, and textual hidden representations
are jointly updated across layers, whereas only the audio latent is evolved
by the flow solver.

\subsubsection{Speaker-Conditioned Audio Refinement}
\label{subsubsec:mmdit_refinement}
After $N_J=12$ joint blocks, the audio representation passes through $N_A=6$ audio-only
blocks with self-attention and feed-forward transformations. A frozen
speaker-verification encoder~\cite{wang2023campplus} provides an
explicit identity embedding, extracted from the complete training utterance
before temporal masking or from the supplied reference utterance at
inference. This embedding is $\ell_2$-normalized, projected, and added to
the flow-time embedding to modulate the audio-only blocks. The preceding
joint blocks remain conditioned on flow time alone. This identity cue
supplements rather than replaces the reference-latent context. Finally, a
time-conditioned normalization and linear projection predict the velocity
field with the same shape as $\mathbf{z}_t$.

\subsubsection{Training Objective}
\label{subsubsec:mmdit_alignment}
We evaluate the flow matching loss only over the masked generation region
and apply auxiliary connectionist temporal classification (CTC)
supervision~\cite{graves2006ctc,choi2025aligndit} to selected audio-stream
hidden states. The overall objective is
\begin{equation}
\mathcal{L}=\mathcal{L}_{\mathrm{CFM}}^{\mathrm{masked}}
+\lambda_{\mathrm{CTC}}\mathcal{L}_{\mathrm{CTC}},
\label{eq:mmdit_training_objective}
\end{equation}
where $\mathcal{L}_{\mathrm{CTC}}$ is averaged over the selected prediction
heads and $\lambda_{\mathrm{CTC}}$ controls its contribution. Transcript
supervision encourages linguistic consistency without externally supplied
phoneme durations. The following subsection describes the temporal bias
used to guide audiovisual interaction.

% Retain the following details in Implementation Details, rather than
% expanding the backbone discussion: hidden width 768, 12 attention heads,
% FFN expansion ratio 2, and 192-dimensional speaker embeddings.
